% Part: proof-theory
% Chapter: sequent-calculus
% Section: invertibility

\documentclass[../../../include/open-logic-section]{subfiles}

\begin{document}

\olfileid{pt}{seq}{inv}

\olsection{وارون‌پذیری قواعد}

\begin{defn}
قاعدهٔ~$R$،
\[
\AxiomC{$S_1$}
\AxiomC{$\dots$}
\AxiomC{$S_n$}
\RightLabel{$R$}
\TrinaryInfC{$S$}
\DisplayProof
\]
در یک دستگاه \emph{وارون‌پذیر} است هرگاه از~$\Proves S$ بتوان برای هر
$i = 1$، \dots،~$n$ نتیجه گرفت که~$\Proves S_i$. این قاعده در دستگاه
\emph{وارون‌پذیر با حفظ !!{height}} (یا
\emph{!!{hp}-وارون‌پذیر}) است هرگاه از~$\Proves[n] S$ بتوان برای هر
$i = 1$، \dots،~$n$ نتیجه گرفت که~$\Proves[n] S_i$.
\end{defn}

\begin{lem}\ollabel{lem:G3c-invert}
قواعد~\Log{G3c} برای~$\lnot$، $\land$، $\lor$ و~$\lif$ با حفظ
!!{height} وارون‌پذیرند؛ یعنی:
\begin{enumerate}
  \item \RightR{\lnot}: اگر~$\Log{G3c} \Proves[n] \Gamma \Sequent
  \Delta, \lnot !A$، آنگاه~$\Log{G3c} \Proves[n] !A, \Gamma \Sequent
  \Delta$.
  \item \LeftR{\lnot}: اگر~$\Log{G3c} \Proves[n] \Gamma, \lnot !A
  \Sequent \Delta$، آنگاه~$\Log{G3c} \Proves[n] \Gamma \Sequent \Delta,
  !A$.
  \item \RightR{\land}: اگر~$\Log{G3c} \Proves[n] \Gamma \Sequent
  \Delta, !A \land !B$، آنگاه~$\Log{G3c} \Proves[n] \Gamma \Sequent
  \Delta, !A$ و~$\Log{G3c} \Proves[n] \Gamma \Sequent
  \Delta, !B$.
  \item \LeftR{\land}: اگر~$\Log{G3c} \Proves[n] \Gamma, !A \land !B
  \Sequent \Delta$، آنگاه~$\Log{G3c} \Proves[n] !A, !B, \Gamma \Sequent
  \Delta$.
  \item \RightR{\lor}: اگر~$\Log{G3c} \Proves[n] \Gamma \Sequent
  \Delta, !A \lor !B$، آنگاه~$\Log{G3c} \Proves[n] \Gamma \Sequent
  \Delta, !A, !B$.
  \item \LeftR{\lor}: اگر~$\Log{G3c} \Proves[n] !A \lor !B, \Gamma
  \Sequent \Delta$، آنگاه~$\Log{G3c} \Proves[n] !A, \Gamma \Sequent
  \Delta$ و~$\Log{G3c} \Proves[n] !B, \Gamma \Sequent
  \Delta$.
  \item \RightR{\lif}: اگر~$\Log{G3c} \Proves[n] \Gamma \Sequent
  \Delta, !A \lif !B$، آنگاه~$\Log{G3c} \Proves[n] !A, \Gamma \Sequent
  \Delta, !B$.
  \item \LeftR{\lif}: اگر~$\Log{G3c} \Proves[n] !A \lif !B, \Gamma
  \Sequent \Delta$، آنگاه~$\Log{G3c} \Proves[n] \Gamma \Sequent \Delta,
  !A$ و~$\Log{G3c} \Proves[n] !B, \Gamma \Sequent \Delta$.
\end{enumerate}
\end{lem}

\begin{proof}
باید نشان دهیم برای هر قاعدهٔ~$R$ در~\Log{G3c}، اگر !!a{proof}~$\pi$
از نتیجهٔ~$S$ قاعدهٔ~$R$ وجود داشته باشد، !!{proof}s~$\pi_i$ از مقدمات
$S_i$ نیز وجود دارند، به‌طوری‌که
$\pheight{\pi_i} \le \pheight{\pi}$. قاعدهٔ~\LeftR{\land} را در نظر
بگیرید: فرض کنید !!a{proof}~$\pi$ با دنبالهٔ پایانی به صورت
$!A \land !B, \Gamma \Sequent \Delta$ وجود دارد. باید نشان دهیم
!!a{proof}~$\pi'$ از~$!A, !B, \Gamma \Sequent \Delta$ وجود دارد که
$\pheight{\pi'} \le \pheight{\pi}$. بر~$n = \pheight{\pi}$ استقرا
می‌کنیم.

اگر~$n = 0$، $\pi$ تنها از یک اصل~$!A \land !B, \Gamma
\Sequent \Delta$ تشکیل می‌شود. اصول~\Log{G3c} یا از صورت
$!C, \Gamma \Sequent \Delta, !C$ با~$!C$ اتمی‌اند، یا اصلِ دارای
$\lfalse$ در سمت چپ‌اند. در صورت نخست،~$!C$ نمی‌تواند~$!A \land !B$
باشد؛ پس یک~$!C$ اتمی !!a{element} از هر دو~$\Gamma$ و~$\Delta$ است و
$!A, !B, \Gamma \Sequent \Delta$ نیز اصل است. در صورت دوم نیز افزودن
$!A$ و~$!B$ آن اصل را حفظ می‌کند. بنابراین !!{proof}~$\pi'$
را یافته‌ایم که آن هم !!{height}~$0$ دارد.

اگر~$n>0$، $\pi$ به استنتاجی ختم می‌شود. باید ادعا را برای~$n$ ثابت
کنیم، اما می‌توانیم در مورد هر !!{proof} با ارتفاع کمتر از~$n$، حتی با
بافتی متفاوت، آن را فرض کنیم. به بیان دیگر، برای هر~$k<n$ و هر~$\Pi$
و~$\Lambda$، اگر~$!A \land !B, \Pi \Sequent \Lambda$ !!a{proof} با
!!{height}~$k$ داشته باشد، آنگاه~$!A, !B, \Pi \Sequent \Lambda$
!!a{proof} با~!!{height}~$\le k$ دارد.

دو حالت داریم: یا~$!A \land !B$ در سمت چپ اصلی است یا نیست. اگر اصلی
باشد، استنتاج پایانی~\LeftR{\land} است و زیر-!!{proof} بی‌واسطهٔ
$\pi_1$ به~$!A, !B, \Gamma \Sequent \Delta$ ختم می‌شود. در این حالت
نتیجه برقرار است، زیرا~$\pheight{\pi_1} < \pheight{\pi}$.

اگر~$!A \land !B$ اصلی نباشد، !!{formula} دیگری اصلی است و
$!A \land !B$ در بافتِ استنتاج پایانی قرار دارد. در این صورت فرض استقرا
بر زیر-!!{proof}s بی‌واسطهٔ قاعدهٔ پایانی~$R$ اعمال می‌شود و یک
!!a{proof}~$\pi_1'$ یا دو~!!{proof}s~$\pi_1'$ و~$\pi_2'$ می‌دهد که
!!{height} آن‌ها از ارتفاع زیر-!!{proof}s بی‌واسطهٔ~$\pi$ بیشتر نیست.
دنباله‌های پایانی~$\pi_1'$ و~$\pi_2'$ به‌جای~$!A \land !B$ در بافت چپ،
$!A, !B$ دارند. با اعمال قاعدهٔ~$R$، برهان~$\pi'$ را به دست
می‌آوریم. برای نمونه، فرض کنید~$\pi$ به~\LeftR{\lif} ختم شود:
\[
\AxiomC{}
\RightLabel{$\pi_1$}
\Deduce$!A \land !B, \Gamma' \fCenter \Delta, !C$
\AxiomC{}
\RightLabel{$\pi_2$}
\Deduce$!A \land !B, \Gamma', !D \fCenter \Delta$
\RightLabel{\LeftR{\lif}}
\BinaryInf$!A \land !B, \Gamma', !C \lif !D \fCenter \Delta$
\DisplayProof
\]
در اینجا~$!C \lif !D$ در سمت چپ اصلی است و بافت چپ
$!A \land !B, \Gamma'$ است؛ یعنی~$\Gamma = \Gamma', !C \lif !D$.
فرض استقرا به‌ترتیب !!{proof}s~$\pi_1'$ و~$\pi_2'$ را از
$!A, !B, \Gamma' \Sequent \Delta, !C$ و
$!A, !B, \Gamma', !D \Sequent \Delta$ می‌دهد. این دو دنباله بار دیگر
مقدمات قاعدهٔ~\LeftR{\lif} می‌شوند تا~$\pi'$ به دست آید:
\[
\AxiomC{}
\RightLabel{$\pi_1'$}
\Deduce$!A, !B, \Gamma' \fCenter \Delta, !C$
\AxiomC{}
\RightLabel{$\pi_2'$}
\Deduce$!A, !B, \Gamma', !D \fCenter \Delta$
\RightLabel{\LeftR{\lif}}
\BinaryInf$!A, !B, \Gamma', !C \lif !D \fCenter \Delta$
\DisplayProof
\]
داریم
\begin{align*}
\pheight{\pi} & = \max(\pheight{\pi_1}, \pheight{\pi_2}) + 1\\
\pheight{\pi'} & = \max(\pheight{\pi_1'}, \pheight{\pi_2'}) + 1
\intertext{طبق تعریف، و}
\pheight{\pi_1'} & \le \pheight{\pi_1}\\
\pheight{\pi_2'} & \le \pheight{\pi_2}
\intertext{طبق فرض استقرا. پس}
\pheight{\pi'} & \le \pheight{\pi}.
\end{align*}
\end{proof}

\begin{prob}
برهان~\olref{lem:G3c-invert} را برای~\RightR{\lif} ارائه کنید.
\end{prob}

\begin{lem}\ollabel{lem:invert-quant}
قواعد~\RightR{\lforall} و~\LeftR{\lexists} به معنای زیر با حفظ
!!{height} وارون‌پذیرند:
\begin{enumerate}
  \item اگر~$\Log{G3c} \Proves[n] \Gamma \Sequent \Delta,
  \lforall[x][!A(x)]$، آنگاه~$\Log{G3c} \Proves[n] \Gamma \Sequent
  \Delta, !A(c)$؛ و
  \item اگر~$\Log{G3c} \Proves[n] \lexists[x][!A(x)], \Gamma \Sequent
  \Delta$، آنگاه~$\Log{G3c} \Proves[n] !A(c), \Gamma \Sequent \Delta$،
\end{enumerate}
برای هر~$c$ که به‌ترتیب در~$\Gamma$، $\Delta$ و فرمول کمیت‌گذاری‌شدهٔ
$\lforall[x][!A(x)]$ یا~$\lexists[x][!A(x)]$ رخ ندهد.
\end{lem}

\begin{proof}
  بند~(1) را نشان می‌دهیم و بند~(2) را به‌عنوان تمرین واگذار می‌کنیم.
  استدلال شبیه برهان~\olref{lem:G3c-invert} است. در گام استقرا باز دو
  حالت داریم. اگر~$\lforall[x][!A(x)]$ اصلی باشد، کار سریع است: مقدمهٔ
  آخرین استنتاج~$\RightR{\lforall}$ به صورت مطلوب
  $\Gamma \Sequent \Delta, !A(a)$ است و زیر-!!{proof} بی‌واسطهٔ~$\pi_1$
  که به آن ختم می‌شود،~$\pheight{\pi_1} < \pheight{\pi}$ دارد. افزون
  بر این، شرط متغیر ویژه می‌گوید~$a$ در~$\Gamma$، $\Delta$ یا
  $\lforall[x][!A(x)]$ رخ نمی‌دهد. چون می‌توان فرض کرد~$\pi_1$ منظم
  است، بنا بر~\olref[qua]{lem:subst-regular}،
  $\Subst{\pi_1}{c}{a}$ !!a{proof} با همان !!{height} از
  $\Gamma \Sequent \Delta, !A(c)$ است.

  در حالت دوم،~$\lforall[x][!A(x)]$ در استنتاج پایانی اصلی نیست و
  فرمول دیگری اصلی است. بار دیگر یک حالت را به‌عنوان نمونه بررسی می‌کنیم.
  فرض کنید استنتاج پایانی~\RightR{\land} باشد. !!{formula}ی در~$\Delta$ به
  صورت~$!B \land !C$ و اصلی است. !!{proof}~$\pi$ چنین پایان می‌یابد:
  \[
\AxiomC{}
\RightLabel{$\pi_1$}
\Deduce$\Gamma \fCenter \Delta', \lforall[x][!A(x)], !B$
\AxiomC{}
\RightLabel{$\pi_2$}
\Deduce$\Gamma \fCenter \Delta', \lforall[x][!A(x)], !C$
\RightLabel{\RightR{\land}}
\BinaryInf$\Gamma \fCenter \Delta', \lforall[x][!A(x)], !B \land !C$
\DisplayProof
\]
که در آن~$\Delta = \Delta', !B \land !C$. !!a{constant}~$c$ را چنان
برگزینید
که در~$\Gamma$، $\Delta$ و~$\lforall[x][!A(x)]$ رخ ندهد؛ پس به‌روشنی
در~$\Delta', !B$ و~$\Delta', !C$ نیز رخ نمی‌دهد. بنابراین فرض استقرا
بر~$\pi_1$ و~$\pi_2$ اعمال می‌شود؛ !!{proof}s~$\pi_1'$ و~$\pi_2'$
داریم که
$\pheight{\pi_1'} \le \pheight{\pi_1}$ و
$\pheight{\pi_2'} \le \pheight{\pi_2}$. !!{proof} مطلوب~$\pi'$ از
$\Gamma \Sequent \Delta, !A(c)$ چنین است:
\[
\AxiomC{}
\RightLabel{$\pi_1'$}
\Deduce$\Gamma \fCenter \Delta', !A(c), !B$
\AxiomC{}
\RightLabel{$\pi_2'$}
\Deduce$\Gamma \fCenter \Delta', !A(c), !C$
\RightLabel{\RightR{\land}}
\BinaryInf$\Gamma \fCenter \Delta', !A(c), !B \land !C$
\DisplayProof
\]
و
$\pheight{\pi'} = \max(\pheight{\pi_1'}, \pheight{\pi_2'})+1
\le \pheight{\pi}$.
\end{proof}

\Olref{lem:G3c-invert} در برهان قضیهٔ حذف برش نقش مهمی دارد، اما برای
نشان‌دادن نتیجهٔ زیر نیز به کار می‌رود:

\begin{prop}\ollabel{prop:G3c-cont-adm}
قواعد انقباض~\LeftR{\Contraction} و~\RightR{\Contraction} در~\Log{G1c}،
در~\Log{G3c} با حفظ !!{height} پذیرفتنی‌اند.
\end{prop}

\begin{proof}
استدلال~\RightR{\Contraction} و~\LeftR{\Contraction} را هم‌زمان
می‌آوریم. فرض کنید~$\pi$ در~\Log{G3c} !!a{proof} از
$\Gamma \Sequent \Delta, !A, !A$ یا
$!A, !A, \Gamma \Sequent \Delta$ باشد. باید نشان دهیم !!{proof}
~\Log{G3c} به نام~$\pi'$ به‌ترتیب از
$\Gamma \Sequent \Delta, !A$ یا~$!A, \Gamma \Sequent \Delta$ وجود دارد
که~$\pheight{\pi'} \le \pheight{\pi}$. بر
$n = \pheight{\pi}$ استقرا می‌کنیم.

اگر~$n=0$، $\pi$ فقط یک اصل است. اگر
$\Gamma \Sequent \Delta, !A, !A$ اصل~\Log{G3c} باشد،
$\Gamma \Sequent \Delta, !A$ نیز اصل است: یا یک~$!C$ اتمی
!!a{element} از هر دو~$\Gamma$ و~$\Delta$ است، یا~$!A$ اتمی
و~!!a{element} از~$\Gamma$ است،
یا مقدم دارای~$\lfalse$ است. همین استدلال برای دنبالهٔ پایانی
$!A, !A, \Gamma \Sequent \Delta$ نیز برقرار است.

اگر~$n>0$، $\pi$ به قاعده‌ای چون~$R$ ختم می‌شود. اگر~$!A$ در این
استنتاج پایانی اصلی نباشد، همانند برهان~\olref{lem:G3c-invert}،
!!a{proof}~$\pi'$ را با اعمال فرض استقرا بر زیر-!!{proof}s بی‌واسطهٔ
$\pi$ و سپس اعمال همان قاعدهٔ~$R$ بر !!{proof}(s) حاصل به دست می‌آوریم.
فرض استقرا این است که اگر~$k<n$ و !!a{proof} با !!{height}~$k$ از
$\Pi \Sequent \Lambda, !D, !D$ یا
$!D, !D, \Pi \Sequent \Lambda$ وجود داشته باشد، آنگاه به‌ترتیب
!!a{proof} با~!!{height}~$\le k$ از
$\Pi \Sequent \Lambda, !D$ یا~$!D, \Pi \Sequent \Lambda$ وجود دارد.
(توجه کنید فرض استقرا حتی وقتی بافت‌ها یا !!{formula} منقبض‌شونده با
$\Gamma$، $\Delta$ یا~$!A$ متفاوت‌اند نیز اعمال می‌شود.)

برای نمونه، فرض کنید~$R$ همان~$\RightR{\land}$ باشد و~$\pi$ چنین پایان
یابد:
\[
\AxiomC{}
\RightLabel{$\pi_1$}
\Deduce$\Gamma \fCenter \Delta', !B, !A, !A$
\AxiomC{}
\RightLabel{$\pi_2$}
\Deduce$\Gamma \fCenter \Delta', !C, !A, !A$
\RightLabel{\RightR{\land}}
\BinaryInf$\Gamma \fCenter \Delta', !B \land !C, !A, !A$
\DisplayProof
\]
که در آن~$\Delta = \Delta', !B \land !C$. فرض استقرا !!{proof}s
$\pi_1'$ و~$\pi_2'$ را به‌ترتیب از
$\Gamma \fCenter \Delta', !B, !A$ و
$\Gamma \fCenter \Delta', !C, !A$ می‌دهد، با
$\pheight{\pi_1'} \le \pheight{\pi_1}$ و
$\pheight{\pi_2'} \le \pheight{\pi_2}$. !!{proof}~$\pi'$ چنین است:
\[
\AxiomC{}
\RightLabel{$\pi_1'$}
\Deduce$\Gamma \fCenter \Delta', !B, !A$
\AxiomC{}
\RightLabel{$\pi_2'$}
\Deduce$\Gamma \fCenter \Delta', !C, !A$
\RightLabel{\RightR{\land}}
\BinaryInf$\Gamma \fCenter \Delta', !B \land !C, !A$
\DisplayProof
\]

اگر~$!A$ در استنتاج پایانی اصلی باشد، لم وارون‌سازی
(~\olref{lem:G3c-invert}) را بر زیر-!!{proof}s بی‌واسطهٔ~$\pi$ اعمال
می‌کنیم، سپس فرض استقرا و بعد دوباره قاعدهٔ~$R$ را. برای نمونه، فرض کنید
$\pi$ دنبالهٔ~$\Gamma \Sequent \Delta, !A, !A$ را اثبات کند،
$!A \ident (!B \land !C)$ باشد و در استنتاج پایانی اصلی باشد. آنگاه
$\pi$ چنین است:
\[
\AxiomC{}
\RightLabel{$\pi_1$}
\Deduce$\Gamma \fCenter \Delta, !B \land !C, !B$
\AxiomC{}
\RightLabel{$\pi_2$}
\Deduce$\Gamma \fCenter \Delta, !B \land !C, !C$
\RightLabel{\RightR{\land}}
\BinaryInf$\Gamma \fCenter \Delta, !B \land !C, !B \land !C$
\DisplayProof
\]
با اعمال~\olref{lem:G3c-invert} بر~$\pi_1$، !!a{proof}~$\pi_1'$ از
$\Gamma \Sequent \Delta, !B, !B$ می‌گیریم. (این لم !!a{proof}ی از
$\Gamma \Sequent \Delta, !C, !B$ نیز می‌دهد که به آن نیاز نداریم.)
به همین ترتیب، اعمال~\olref{lem:G3c-invert} بر~$\pi_2$،
!!a{proof}~$\pi_2'$ از~$\Gamma \Sequent \Delta, !C, !C$ می‌دهد. چون
وارون‌سازی~\RightR{\land} !!{height} را حفظ می‌کند،
$\pheight{\pi_1'} \le \pheight{\pi_1} < \pheight{\pi}$ و
$\pheight{\pi_2'} \le \pheight{\pi_2} < \pheight{\pi}$. پس فرض استقرا
بر~$\pi_1'$ و~$\pi_2'$ اعمال می‌شود: !!{proof}s~$\pi_1''$ و~$\pi_2''$
به‌ترتیب از~$\Gamma \Sequent \Delta, !B$ و
$\Gamma \Sequent \Delta, !C$ وجود دارند، با
$\pheight{\pi_1''} \le \pheight{\pi_1'} \le \pheight{\pi_1}$ و
$\pheight{\pi_2''} \le \pheight{\pi_2'} \le \pheight{\pi_2}$.
آنگاه !!{proof}~$\pi'$ چنین است:
\[
\AxiomC{}
\RightLabel{$\pi_1''$}
\Deduce$\Gamma \fCenter \Delta, !B$
\AxiomC{}
\RightLabel{$\pi_2''$}
\Deduce$\Gamma \fCenter \Delta, !C$
\RightLabel{\RightR{\land}}
\BinaryInf$\Gamma \fCenter \Delta, !B \land !C$
\DisplayProof
\]
و~$\pheight{\pi'} \le \pheight{\pi}$.

در حالت‌هایی که !!{formula} کمیت‌گذار و اصلی است باید دقت ویژه‌ای به خرج
دهیم.

فرض کنید~$!A \ident \lforall[x][!B(x)]$ در سمت راست اصلی باشد. آنگاه
$\pi$ چنین پایان می‌یابد:
\[
\AxiomC{}
\RightLabel{$\pi_1$}
\Deduce$\Gamma \fCenter \Delta, \lforall[x][!B(x)], !B(a)$
\RightLabel{\RightR{\lforall}}
\UnaryInf$\Gamma \fCenter \Delta, \lforall[x][!B(x)], \lforall[x][!B(x)]$
\DisplayProof
\]
بنا بر~\olref{lem:invert-quant}، برای هر~$c$ که در~$\Gamma$، $\Delta$،
$\lforall[x][!B(x)]$ و~$!B(a)$ رخ ندهد، !!a{proof}~$\pi_1'$ از
$\Gamma \fCenter \Delta, !B(c), !B(a)$ با !!{height}~$\le$
$\pheight{\pi_1}$ وجود دارد. می‌توان فرض کرد~$\pi_1'$ منظم است؛ پس بنا
بر~\olref[qua]{lem:subst-regular}، !!{proof}
$\Subst{\pi_1'}{a}{c}$، !!a{proof} از
$\Gamma \fCenter \Delta, !B(a), !B(a)$ با
!!{height}~$\le \pheight{\pi_1}$ است. اکنون فرض استقرا اعمال می‌شود و
!!a{proof}~$\pi_1''$ از~$\Gamma \fCenter \Delta, !B(a)$ می‌دهد. با
قاعدهٔ~\RightR{\lforall}، $\pi'$ را می‌گیریم:
\[
\AxiomC{}
\RightLabel{$\pi_1''$}
\Deduce$\Gamma \fCenter \Delta, !B(a)$
\RightLabel{\RightR{\lforall}}
\UnaryInf$\Gamma \fCenter \Delta, \lforall[x][!B(x)]$
\DisplayProof
\]
و~$\pheight{\pi'} \le \pheight{\pi}$.

اکنون فرض کنید~$!A \ident \lexists[x][!B(x)]$ و در سمت راست اصلی باشد.
آنگاه~$\pi$ چنین است:
\[
\AxiomC{}
\RightLabel{$\pi_1$}
\Deduce$\Gamma \fCenter \Delta, \lexists[x][!B(x)], \lexists[x][!B(x)], !B(t)$
\RightLabel{\RightR{\lexists}}
\UnaryInf$\Gamma \fCenter \Delta, \lexists[x][!B(x)], \lexists[x][!B(x)]$
\DisplayProof
\]
فرض استقرا بر~$\pi_1$ اعمال می‌شود و !!a{proof}~$\pi_1'$ از
$\Gamma \Sequent \Delta, \lexists[x][!B(x)], !B(t)$ وجود دارد.
(اینجا به هیچ لم وارون‌سازی نیاز نیست.) آنگاه !!{proof}~$\pi'$ چنین است:
\[
\AxiomC{}
\RightLabel{$\pi_1'$}
\Deduce$\Gamma \fCenter \Delta, \lexists[x][!B(x)], !B(t)$
\RightLabel{\RightR{\lexists}}
\UnaryInf$\Gamma \fCenter \Delta, \lexists[x][!B(x)]$
\DisplayProof
\]
و~$\pheight{\pi'} \le \pheight{\pi}$.
\end{proof}

\begin{prob}
برهان~\olref{prop:G3c-cont-adm} را برای~\LeftR{\Contraction} طرح کنید.
حالت‌هایی را شرح دهید که~$!A \ident (!B \land !C)$ و
$!A \ident (!B \lif !C)$.
\end{prob}

\begin{prob}
برهان~\olref{prop:G3c-cont-adm} را برای~\LeftR{\Contraction} در حالت‌های
کمیت‌گذارِ باقی‌مانده طرح کنید؛ یعنی وقتی~$!A \ident
\lforall[x][!B(x)]$ و~$!A \ident \lexists[x][!B(x)]$.
\end{prob}


\end{document}
